%!TEX program = xelatex
\def\numeromodulo{1}
\def\numerolista{1}
\documentclass[12pt]{exam}

\def\ano{2023}
\def\semestre{1}
\def\disciplina{Introdução à \'Algebra Linear}
\def\turma{3}

\usepackage{caption}
\usepackage{amssymb}
\usepackage{amsmath,amsfonts,amsthm,amstext}
\usepackage[utf8]{inputenc} % Required for inputting international characters
\usepackage[T1]{fontenc}
\usepackage[brazil]{babel}
\usepackage{graphicx}
\graphicspath{{../figuras_geogebra/}{../Pictures/}}
\usepackage{enumitem}
\usepackage{multicol}
\usepackage{answers}
\usepackage[svgnames]{xcolor}
\usepackage{tikz}
\usepackage{ifthen}
\usetikzlibrary{lindenmayersystems}
\usetikzlibrary[shadings]
\usepackage{tikz-3dplot}
\usepackage{tcolorbox}
%-------------amatrix
% Augmented matrix.  Usage (note the argument does not count the aug col):
% \begin{amatrix}{2}
%   1  2  3 \\  4  5  6
% \end{amatrix}
\newenvironment{amatrix}[1]{%
  \left(\begin{array}{@{}*{#1}{r}|r@{}}
}{%
  \end{array}\right)
}

\Newassociation{solucao}{Solution}{ans}
\newtheorem{exercicio}{}

\setlength{\topmargin}{-1.0in}
\setlength{\oddsidemargin}{0in}
\setlength{\textheight}{10.1in}
\setlength{\textwidth}{6.5in}
\setlength{\baselineskip}{12mm}

\newcounter{exercicios}
\setcounter{exercicios}{0}
\newcommand{\questao}{
\addtocounter{exercicios}{1}
\noindent{\bf Exerc{\'\i}cio \Roman{exercicios}: }}

\newcommand{\resp}[1]{
\noindent{\bf Exerc{\'\i}cio #1: }}

\extraheadheight{0.7in}
\firstpageheadrule
\runningheadrule
\lhead{
        \begin{minipage}[c]{1.7cm}
        \includegraphics[width=1.7cm]{unb.pdf}
        \end{minipage}%
        \hspace{0pt}
        \begin{minipage}[c]{4in}
          {Universidade de Bras{\'\i}lia} --
          {Departamento de Matem{\'a}tica}
\end{minipage}
\vspace*{-0.8cm}
}
% \chead{Universidade de Bras{\'\i}lia - Departamento de Matem{\'a}tica}
% \rhead{}
% \vspace*{-2cm}

\extrafootheight{.5in}
\footrule
\lfoot{\disciplina\ - \semestre$^o$/\ano\ - Módulo \numeromodulo}
\cfoot{}
\rfoot{P\'agina \thepage\ de \numpages}
\newcommand{\nsub}{\varsubsetneq}
\newcommand{\vaz}{\emptyset}
\newcommand{\im}{{\rm Im\,}}
\newcommand{\sub}{\subseteq}
\renewcommand{\sin}{{\rm sen\,}}
\newcommand{\n}{\mathbb{N}}
\newcommand{\z}{\mathbb{Z}}
\newcommand{\rac}{\mathbb{Q}}
\newcommand{\real}{\mathbb{R}}
\newcommand{\complex}{\mathbb{C}}
\newcommand{\cp}[1]{\mathbb{#1}}
\newcommand{\ch}{\mbox{\textrm{car\,}}\nobreak}
\newcommand{\vesp}[1]{\vspace{ #1  cm}}
\newcommand{\compcent}[1]{\vcenter{\hbox{$#1\circ$}}}
\newcommand{\comp}{\mathbin{\mathchoice
{\compcent\scriptstyle}{\compcent\scriptstyle}
{\compcent\scriptscriptstyle}{\compcent\scriptscriptstyle}}}

\newcommand{\integer}{\mathbb{Z}}
\begin{document}

  \Opensolutionfile{ans}[ans1]
  \begin{center}
    {\Large\bf \disciplina\ - Turma \turma\ -- \semestre$^{o}$/\ano} \\ \vspace{9pt} {\large\bf
        $\numerolista^a$ Lista de Exerc{\'\i}cios -- Módulo \numeromodulo}\\ \vspace{9pt} Prof. Jos{\'e} Ant{\^o}nio O. Freitas
  \end{center}
  \hrule

\vesp{.6}

\begin{exercicio}
  \label{operacoes_matrizes} Sejam
    \begin{align*}
        A = \begin{pmatrix}1 & 2 & \phantom{-} 3\\2 & 1 & -1\end{pmatrix},
        B = \begin{pmatrix}-2 & 0 & 1\\\phantom{-} 3 & 0 & 1\end{pmatrix},
        C = \begin{pmatrix}-1\\ \phantom{-} 2\\\phantom{-} 4\end{pmatrix},
        D = \begin{pmatrix}-2\\ \phantom{-} 3\\\phantom{-} 2\end{pmatrix},
        E = \begin{pmatrix}2 & -1\end{pmatrix},
        F = \begin{pmatrix}4 & 3\end{pmatrix}.
    \end{align*}
    Encontre:
    \begin{multicols}{3}
        \begin{enumerate}[label={\alph*})]
            \item $A + B$
            \item $3B$
            \item $2A - B$
            \item $3(A - \dfrac{1}{2}B)$
            \item $A + 2B$
            \item $2C + D$
            \item $\dfrac{2}{3}D - C$
            \item $3E + 4F$
            \item $E - F$
            \item $F - 2E$
            \item $AC$
            \item $BC$
            \item $CD$
            \item $CB$
            \item $DA$
            \item $DB$
            \item $-A$
            \item $-2F$
        \end{enumerate}
    \end{multicols}
  \begin{solucao}
    \begin{enumerate}[label={\alph*})]
        \item $A + B = \begin{pmatrix} -1 & 2 & 4\\5 & 1 & 0\end{pmatrix}$
        \item $3B = \begin{pmatrix} -6 & 0 & 3\\9 & 0 & 3\end{pmatrix}$
        \item $2A - B = \begin{pmatrix} 4 & 4 & 5\\1 & 2 & -3\end{pmatrix}$
        \item $3(A - \dfrac{1}{2}B) = \begin{pmatrix} 0 & 6 & 21/2\\21/2 & 3 & -3/2\end{pmatrix}$
        \item $A + 2B = \begin{pmatrix} -3 & 2 & 5\\8 & 1 & 1\end{pmatrix}$
        \item $2C + D = \begin{pmatrix} -4\\ 7\\ 10\end{pmatrix}$
        \item $\dfrac{2}{3}D - C = \begin{pmatrix} -1/3\\ 0\\ -8/3\end{pmatrix}$
        \item $3E + 4F = \begin{pmatrix} 22 & 9\end{pmatrix}$
        \item $E - F = \begin{pmatrix} -2 & -4\end{pmatrix}$
        \item $F - 2E = \begin{pmatrix} 0 & 5\end{pmatrix}$
        \item $AC = \begin{pmatrix} 15\\ -4\end{pmatrix}$
        \item $BC = \begin{pmatrix} 6\\1\end{pmatrix}$
        \item O produto $CD$ não está definido.
        \item O produto $CB$ não está definido.
        \item O produto $DA$ não está definido.
        \item O produto $DB$ não está definido.
        \item $-A = \begin{pmatrix} -1 & -2 & 3\\-2 & -1 & 1\end{pmatrix}$
        \item $2F = \begin{pmatrix} -8 & -6\end{pmatrix}$
    \end{enumerate}
  \end{solucao}
\end{exercicio}

\begin{exercicio}
    Nos itens abaixo são dadas duas matrizes $A$ e $B$ e dois números $c$ e $d$. Calcule a matriz $cA + dB$.
    \begin{enumerate}[label={\alph*})]
        \item $A = \begin{pmatrix}3 & -5\\2 & \phantom{-} 7\end{pmatrix}$, $B = \begin{pmatrix}-1 & \phantom{-} 0\\\phantom{-} 3 & -4\end{pmatrix}$, $c = 3$, $d = 4$
        \item $A = \begin{pmatrix}\phantom{-} 2 & 0 & -3\\-1 & 5 & \phantom{-} 6\end{pmatrix}$, $B = \begin{pmatrix}-2 & 3 & 1\\\phantom{-} 7 & 1 & 5\end{pmatrix}$, $c = 5$, $d = -3$
        \item $A = \begin{pmatrix}5 & \phantom{-} 0\\0 & \phantom{-} 7\\3 & -1\end{pmatrix}$, $B = \begin{pmatrix}-4 & 5\\\phantom{-} 3 & 2\\\phantom{-} 7 & 4\end{pmatrix}$, $c = -2$, $d = 4$
        \item $A = \begin{pmatrix}2 & -1 & \phantom{-} 0\\4 & \phantom{-} 0 & -3\\5 & -2 & \phantom{-} 7\end{pmatrix}$, $B = \begin{pmatrix}6 & -3 & -4\\5 & \phantom{-} 2 & -1\\0 & \phantom{-} 7 & \phantom{-} 9\end{pmatrix}$, $c = 7$, $d = 5$
    \end{enumerate}
    \begin{solucao}
        \begin{enumerate}[label={\alph*})]
            \item $3A + 4B = \begin{pmatrix}-13 & -15\\18 & 5\end{pmatrix}$
            \item $5A - 3B = \begin{pmatrix}16 & -9 & -18\\-26 & 22 & 15\end{pmatrix}$
            \item $-2A + 4B = \begin{pmatrix}-26 & 20\\12 & -6\\22 & 18\end{pmatrix}$
            \item $7A + 5B = \begin{pmatrix}44 & -22 & -20\\53 & 10 & 16\\35 & 21 & 94\end{pmatrix}$
        \end{enumerate}
    \end{solucao}
\end{exercicio}

\begin{exercicio}
    \label{igualdade_de_matrizes} Nos itens abaixo encontre os valores de $a$, $b$, $c$ e $d$ de modo que as igualdades sejam verdadeiras:
    \begin{enumerate}[label={\alph*})]
        \item $\begin{bmatrix}\phantom{-} a & 3\\-1 & a + b\end{bmatrix} = \begin{bmatrix}4 & d - 2c\\d - 2c & -2\end{bmatrix}$
        \item $\begin{bmatrix}a - b & b + a\\3d + c & 2d - c\end{bmatrix} = \begin{bmatrix}8 & 1\\7 & 6\end{bmatrix}$.
    \end{enumerate}
    \begin{solucao}
        \begin{enumerate}
            \item Não existem $d$ e $c$ tais que essas matrizes sejam iguais.
            \item $a = 9/2$, $b = -7/2$, $c = -4$, $d = 1$
        \end{enumerate}
    \end{solucao}
\end{exercicio}

\begin{exercicio}
    Suponha que $A$, $B$, $C$, $D$ e $E$ sejam matrizes tais que:
    \begin{align*}
        A &= (a_{ij})_{4 \times 5}
        B = (b_{ij})_{4 \times 5}
        C = (c_{ij})_{5 \times 2}\\
        D &= (d_{ij})_{4 \times 2}
        E = (e_{ij})_{5 \times 4}.
    \end{align*}
    Determine quais das operações a seguir estão definidas. Para aquelas que estiverem definidas, determine a ordem da matriz resultante. Quando não estiver definida, explique por que não está definida.
    \begin{multicols}{2}
        \begin{enumerate}[label={\alph*})]
            \item $AB$

            \item $BC$

            \item $AC + D$

            \item $AE + B$

            \item $AB + B$

            \item $E(A + B)$

            \item $E(AC)$
        \end{enumerate}
    \end{multicols}
    \begin{solucao}
        \begin{enumerate}[label={\alph*})]
            \item Não está definida. Observer as ordens das matrizes.
            \item $BC$ é uma matriz de ordem $4 \times 2$.
            \item $AC + D$ é uma matriz de ordem $4 \times 2$.
            \item Não está definida. Observer as ordens das matrizes.
            \item Não está definida. Observer as ordens das matrizes.
            \item $E(A + B)$ é uma matriz de ordem $5 \times 5$.
            \item $E(AC)$ é uma matriz de ordem $5 \times 2$.
        \end{enumerate}
    \end{solucao}
\end{exercicio}

\begin{exercicio}
    No Exercício \ref{operacoes_matrizes} encontre $A^t$, $B^t$, $C^t$, $D^t$, $E^t$ e $F^t$.
    \begin{solucao}
        $A^t = \begin{pmatrix}1 & 2\\2 & 1\\3 & -1\end{pmatrix}$
        $B^t = \begin{pmatrix}-2 & 3\\0 & 0\\1 & 1\end{pmatrix}$
        $C^t = \begin{pmatrix}-1 & 2 & 4\end{pmatrix}$
        $D^t = \begin{pmatrix}-2 & 3 & 2\end{pmatrix}$
        $E^t = \begin{pmatrix}2\\ -1\end{pmatrix}$
        $F^t = \begin{pmatrix}4\\3\end{pmatrix}$
    \end{solucao}
\end{exercicio}

\begin{exercicio}
    Encontre a transposta de todas as matrizes que aparecem no Exercício \ref{igualdade_de_matrizes}.
\end{exercicio}

\begin{exercicio}
    Considere as matrizes
    \begin{align*}
        A = \begin{pmatrix}\phantom{-} 5 & \phantom{-} 0 & \phantom{-} 6\\-8 & \phantom{-} 0 & \phantom{-} 3\\-2 & \phantom{-} 2 & \phantom{-} 7\\1 & -1 & -5\end{pmatrix},
        B = \begin{pmatrix}1 & 3 & \phantom{-} 2 & \phantom{-} 4\\7 & 1 & -3 & \phantom{-} 4\\0 & 0 & \phantom{-} 3 & -2\end{pmatrix},
        C = \begin{pmatrix}\phantom{-} 1 & 0 & \phantom{-} 2\\\phantom{-} 1 & 0 & \phantom{-} 1\\-2 & 1 & 2\end{pmatrix},
        D = \begin{pmatrix}\phantom{-} 5 & \phantom{-} 0 & 1 & 1\\-3 & \phantom{-} 1 & 4 & 1\\\phantom{-} 2 & \phantom{-} 2 & 0 & 3\\\phantom{-} 1 & -1 & 0 & 0\end{pmatrix}.
    \end{align*}
    Calcule os seguintes produtos:
    \begin{enumerate}[label={\alph*})]
        \item $(AB)D^t$
        \item $3(A^tB^t) + C^t$
        \item $C^tB$
    \end{enumerate}
    \begin{solucao}
        \begin{enumerate}
            \item $(AB)D^t = \begin{pmatrix}66 & 135 & 54 & -5\\-93 & -90 & -162 & 8\\
                55 & -16 & -22 & 14\\ -29 & -7 & 20 & -7\end{pmatrix}$
            \item $3(A^tB^t) + C^t = \begin{pmatrix}-38 & 112 & -26\\6 & -30 & 25\\41 & 13 & 95\end{pmatrix}$
            \item $C^tB = \begin{pmatrix}8 & 3 & -6 & 12\\0 & 0 & 3 & -2\\9 & 5 & 9 & 8\end{pmatrix}$
        \end{enumerate}
    \end{solucao}
\end{exercicio}

\begin{exercicio}
    Para as matrizes do Exercício \ref{operacoes_matrizes} calcule os produtos a seguir, se estiverem definidos.
    \begin{enumerate}[label={\alph*})]
        \item $E^tA$
        \item $(A^t + E)D$
        \item $C^tB$
        \item $B^tE$
        \item $A^tB$
        \item $B^tD$
    \end{enumerate}
    \begin{solucao}
        \begin{enumerate}[label={\alph*})]
            \item $E^tA$ não está definido.
            \item $(A^t + E)D$ não está definido.
            \item $C^tB$ não está definido.
            \item $B^tE$ não está definido.
            \item $A^tB = \begin{pmatrix}4 & 0 & 3\\-1 & 0 & 3\\-9 & 0 & 2\end{pmatrix}$
            \item $B^tD$ não está definido.
        \end{enumerate}
    \end{solucao}
\end{exercicio}

\begin{exercicio}
    Seja
    \[
        A = \begin{bmatrix}2 & x^2\\2x - 1 & 0\end{bmatrix}.
    \]
    Encontre $x$ tal que
    \[
        A^t = A.
    \]
    \begin{solucao}
        $x = 1$
    \end{solucao}
\end{exercicio}

\begin{exercicio}
    Em cada item abaixo, encontre a matriz $A = (a_{ij})$ de ordem $4 \times 4$ satisfazendo a condição dada.
    \begin{enumerate}[label={\alph*})]
        \item $a_{ij} = i + j$
        \item $a_{ij} = i^{j - 1}$
        \item $a_{ij} = \begin{cases}\phantom{-} 1 & \mbox{ se } |i - j| > 1\\-1 & \mbox{ se } |i - j| \le 1\end{cases}$
    \end{enumerate}
\end{exercicio}

\begin{exercicio}
    Em cada item abaixo, encontre uma matriz $A = (a_{ij})$ de ordem $5 \times 5$ satisfazendo a condição dada. Dê um exemplo com números e também escreva a forma geral usando letras ao invés de números.
    \begin{enumerate}[label={\alph*})]
        \item $a_{ij} = 0$ se $i \ne j$
        \item $a_{ij} = 0$ se $i > j$
        \item $a_{ij} = 0$ se $i < j$
        \item $a_{ij} = 0$ se $|i - j| > 1$.
    \end{enumerate}
\end{exercicio}

\begin{exercicio}
    Quais das seguintes matrizes são simétricas?
    \begin{align*}
        A = \begin{pmatrix}2 & -1\\1 & \phantom{-} 2\end{pmatrix},
        B = \begin{pmatrix}\pi & 0\\0 & \sqrt{2}\end{pmatrix},
        C = \begin{pmatrix}-2 & 1 & \phantom{-} 4\\\phantom{-} 1 & 3 & \phantom{-} 1\\\phantom{-} 4 & 3 & -1\end{pmatrix},
        D = \begin{pmatrix}0 & 0 & \phantom{-} 1\\0 & 2 & -1\\1 & 3 & \phantom{-} 4\end{pmatrix}.
    \end{align*}
    \begin{solucao}
        A única matriz simétrica é a matriz $B$.
    \end{solucao}
\end{exercicio}

\begin{exercicio}
    Calcule o \textit{traço} das seguintes matrizes:
    \begin{enumerate}[label={\alph*})]
        \item $A = \begin{pmatrix}2 & \phantom{-} 3 & \phantom{-} 0\\1 & -2 & -3\\3 & -5 & \phantom{-} 2\end{pmatrix}$
        \item $B = \begin{pmatrix}\phantom{-} 5 & \phantom{-} 3 & 0 & \phantom{-} 2 & \phantom{-} 4\\-1 & \phantom{-} 1 & 0 & \phantom{-} 2 & -8\\\phantom{-} 0 & -5 & 2 & \phantom{-} 0 & -1\\\phantom{-} 2 & \phantom{-} 3 & 4 & -3 & -3\end{pmatrix}$
        \item $C = \begin{pmatrix}1 & \phantom{-} 0 & \phantom{-} 0 & 0\\0 & -2 & -3 & 1\\2 & \phantom{-} 4 & -2 & 2\\4 & -1 & \phantom{-} 3 & 0\end{pmatrix}$
    \end{enumerate}
\end{exercicio}

\begin{exercicio}
    Considere a matriz
        \[
            A = \begin{pmatrix}a & b & -5\\c & d & \phantom{-} 0\\5 & 0 & \phantom{-} 0\end{pmatrix}.
        \]
        \begin{enumerate}[label={\alph*})]
            \item Se $A$ é simétrica, encontre o valor de $2a + b + c - 3d$.
            \item Se $A$ é anti-simétrica, encontre o valor de $2a + b + c - 3d$.
        \end{enumerate}
\end{exercicio}

\begin{exercicio}
    Se a matriz
        \[
            A = \begin{pmatrix}1 & 3 + x & 2\\1 - x & 2 & y + 1\\2 & 5 - y & 3\end{pmatrix}
        \]
        é simétrica, qual é o valor de $3x + y$?
\end{exercicio}

\newpage
\Closesolutionfile{ans}
\hrule
\begin{center}
{\large\bf RESPOSTAS}
\end{center}
\hrule
\input{ans1}

\end{document}
